% GENERATED FILE. Edit canonical lessons, then run npm run generate:books.
% canonical-source-hash: fnv1a64:c29ca4b7e0c823c4
% canonical-lessons: PT-C410-curva, PT-C410-lomba, PT-C410-atalho, PT-C410-alcatrao, PT-C410-serpentear

\chapter{Estradas do campo}
\label{ch:pt-estradas-do-campo}
\hlchaptermodality{410}

\begin{chapteropening}
\textbf{By the end of this chapter:} \emph{I can describe country roads: bends, speed bumps, shortcuts, tarmac and roads that wind.}

Complete the last lesson of chapter 410: I can describe country roads: bends, speed bumps, shortcuts, tarmac and roads that wind.
\end{chapteropening}

\section[a curva]{a curva — the bend, the curve}

\label{lesson:PT-C410-curva}



\begin{quote}
Before the new one: say the Portuguese for to crash into, then the Portuguese for the diversion.
\end{quote}

\subsection*{You'll want to know: a curva}
\textbf{a curva} — \textquotedblleft{}the bend, the curve\textquotedblright{}.

\begin{quote}
Depois da curva, a estrada desce até ao rio.
\end{quote}

\begin{quote}
\emph{After the bend, the road goes down to the river.}
\end{quote}

\emph{Uma curva perigosa} — a dangerous bend.

\begin{cousinweb}
\emph{curva} comes from Latin \emph{curvus} 'bent, crooked' (Lewis \& Short); EtymDB has no entry for the Portuguese word.
\end{cousinweb}

\begin{grammarlens}[title={What you've built}]
One more word for telling, explaining and describing at B1.
\end{grammarlens}

\subsection*{Your turn}
\begin{itemize}
\raggedright
  \item \emph{Say it:} \emph{a curva}
  \item \emph{Say it:} \emph{a curva}, once more
  \item \emph{Say it:} a sentence of your own with \emph{a curva}
  \item \emph{Recall:} say the Portuguese for to crash into, then the Portuguese for the diversion, then say \emph{a curva} again
\end{itemize}

\begin{tcolorbox}[breakable,colback=teal!4,colframe=teal!35!black,title={Before you move on}]
What does a curva mean? (\textquotedblleft{}The bend, the curve\textquotedblright{}.) Say it once more.
\end{tcolorbox}
\section[a lomba]{a lomba — the speed bump; the rise (in a road)}

\label{lesson:PT-C410-lomba}



\begin{quote}
Before the new one: say the Portuguese for the diversion, then the Portuguese for the bend.
\end{quote}

\subsection*{You'll want to know: a lomba}
\textbf{a lomba} — \textquotedblleft{}the speed bump; the rise (in a road)\textquotedblright{}.

\begin{quote}
Há lombas à porta da escola.
\end{quote}

\begin{quote}
\emph{There are speed bumps outside the school.}
\end{quote}

Also a rounded hilltop.

\begin{cousinweb}
\emph{lomba} is built on \emph{lombo} 'back, loin', which is inherited from Latin \emph{lumbus} 'loin' (Lewis \& Short): a hump like an animal's back.
\end{cousinweb}

\begin{grammarlens}[title={What you've built}]
One more word for telling, explaining and describing at B1.
\end{grammarlens}

\subsection*{Your turn}
\begin{itemize}
\raggedright
  \item \emph{Say it:} \emph{a lomba}
  \item \emph{Say it:} \emph{a lomba}, once more
  \item \emph{Say it:} a sentence of your own with \emph{a lomba}
  \item \emph{Recall:} say the Portuguese for the diversion, then the Portuguese for the bend, then say \emph{a lomba} again
\end{itemize}

\begin{tcolorbox}[breakable,colback=teal!4,colframe=teal!35!black,title={Before you move on}]
What does a lomba mean? (\textquotedblleft{}The speed bump; the rise (in a road)\textquotedblright{}.) Say it once more.
\end{tcolorbox}
\section[o atalho]{o atalho — the shortcut}

\label{lesson:PT-C410-atalho}



\begin{quote}
Before the new one: say the Portuguese for the bend, then the Portuguese for the speed bump.
\end{quote}

\subsection*{You'll want to know: o atalho}
\textbf{o atalho} — \textquotedblleft{}the shortcut\textquotedblright{}.

\begin{quote}
Conheço um atalho pelo campo.
\end{quote}

\begin{quote}
\emph{I know a shortcut through the fields.}
\end{quote}

\emph{Ir por um atalho} — to take a shortcut.

\begin{cousinweb}
\emph{atalho} is built on \emph{atalhar} 'to cut short', on \emph{talhar} 'to cut', which goes back to Late Latin \emph{taleāre}, from \emph{tālea} 'cutting, stick' (Lewis \& Short).
\end{cousinweb}

\begin{grammarlens}[title={What you've built}]
One more word for telling, explaining and describing at B1.
\end{grammarlens}

\subsection*{Your turn}
\begin{itemize}
\raggedright
  \item \emph{Say it:} \emph{o atalho}
  \item \emph{Say it:} \emph{o atalho}, once more
  \item \emph{Say it:} a sentence of your own with \emph{o atalho}
  \item \emph{Recall:} say the Portuguese for the bend, then the Portuguese for the speed bump, then say \emph{o atalho} again
\end{itemize}

\begin{tcolorbox}[breakable,colback=teal!4,colframe=teal!35!black,title={Before you move on}]
What does o atalho mean? (\textquotedblleft{}The shortcut\textquotedblright{}.) Say it once more.
\end{tcolorbox}
\section[o alcatrão]{o alcatrão — tarmac, tar}

\label{lesson:PT-C410-alcatrao}



\begin{quote}
Before the new one: say the Portuguese for the speed bump, then the Portuguese for the shortcut.
\end{quote}

\subsection*{You'll want to know: o alcatrão}
\textbf{o alcatrão} — \textquotedblleft{}tarmac, tar\textquotedblright{}.

\begin{quote}
A estrada nova tem alcatrão; a velha é de terra.
\end{quote}

\begin{quote}
\emph{The new road is tarmac; the old one is dirt.}
\end{quote}

In Brazil: \emph{o asfalto}.

\begin{cousinweb}
\emph{alcatrão} comes from Arabic \emph{qaṭrān} 'tar' with the article \emph{al-} (EtymDB), like Spanish \emph{alquitrán}.
\end{cousinweb}

\begin{grammarlens}[title={What you've built}]
One more word for telling, explaining and describing at B1.
\end{grammarlens}

\subsection*{Your turn}
\begin{itemize}
\raggedright
  \item \emph{Say it:} \emph{o alcatrão}
  \item \emph{Say it:} \emph{o alcatrão}, once more
  \item \emph{Say it:} a sentence of your own with \emph{o alcatrão}
  \item \emph{Recall:} say the Portuguese for the speed bump, then the Portuguese for the shortcut, then say \emph{o alcatrão} again
\end{itemize}

\begin{tcolorbox}[breakable,colback=teal!4,colframe=teal!35!black,title={Before you move on}]
What does o alcatrão mean? (\textquotedblleft{}Tarmac, tar\textquotedblright{}.) Say it once more.
\end{tcolorbox}
\section[serpentear]{serpentear — to wind, to snake (of a road or river)}

\label{lesson:PT-C410-serpentear}



\begin{quote}
Before the new one: say the Portuguese for the shortcut, then the Portuguese for tarmac.
\end{quote}

\subsection*{You'll want to know: serpentear}
\textbf{serpentear} — \textquotedblleft{}to wind, to snake (of a road or river)\textquotedblright{}.

\begin{quote}
A estrada serpenteia pela montanha.
\end{quote}

\begin{quote}
\emph{The road winds up the mountain.}
\end{quote}

Built on \emph{a serpente}, the snake.

\begin{cousinweb}
\emph{serpentear} is built on \emph{serpente} 'snake', from Latin \emph{serpēns} 'creeping thing', from \emph{serpere} 'to creep' (Lewis \& Short).
\end{cousinweb}

\begin{grammarlens}[title={What you've built}]
One more word for telling, explaining and describing at B1.
\end{grammarlens}

\subsection*{Your turn}
\begin{itemize}
\raggedright
  \item \emph{Say it:} \emph{serpentear}
  \item \emph{Say it:} \emph{serpentear}, once more
  \item \emph{Say it:} a sentence of your own with \emph{serpentear}
  \item \emph{Recall:} say the Portuguese for the shortcut, then the Portuguese for tarmac, then say \emph{serpentear} again
\end{itemize}

\begin{tcolorbox}[breakable,colback=teal!4,colframe=teal!35!black,title={Before you move on}]
What does serpentear mean? (\textquotedblleft{}To wind, to snake (of a road or river)\textquotedblright{}.) Say it once more.
\end{tcolorbox}
